\section{Теоретични основи на \ORMName{} и моделите за съхранение}
\label{sec:theoretical_background}

Теоретичната основа на настоящата разработка не се свежда до една езикова възможност или до един тип база данни. Библиотеката комбинира идеи от ORM теорията, статичното типизиране в C++, compile-time изграждането на междинни представяния и сравнителния анализ на различни модели за съхранение. За да бъде разбираема по-късната архитектура, е необходимо първо да се отдели логическият слой на описване на данните от физическия слой на тяхното backend-специфично материализиране.

В контекста на тази дипломна работа терминът \ORMName{} не трябва да се разбира стеснено само като SQL техника. По-точно той обозначава архитектурен подход, при който моделът на данните, формата на заявките и част от ограниченията върху операциите се описват чрез типовата система на езика, а не чрез runtime конфигурация. Това позволява една и съща логическа структура да бъде пренесена към релационни, документни, key-value, графови и wide-column системи, без backend-специфичните различия да се скриват или подменят.

\subsection{Таксономия на problem domain-а}

Преди разглеждане на отделните компоненти е полезно да се фиксира общата картина на problem domain-а. От една страна стои C++ моделът, в който приложението описва entity типове, полета, връзки и ограничения. От друга страна стоят различни семейства системи за съхранение, всяко със собствен модел на физическа организация, достъп и оптимизация. Между тях е разположен типизираният query IR, който служи като посредник между логиката на приложението и конкретния connector слой.

\begin{figure}[H]
\centering
\begin{tikzpicture}[node distance=1.5cm and 1.4cm]
    \node[ormaccent] (cpp) {\textbf{C++ логически модел}\\entity типове, полета, връзки};
    \node[ormblock, below=of cpp] (ir) {\textbf{Типизиран query IR}\\fields, rules, placeholders};
    \node[ormblock, below left=of ir] (rel) {Релационен\\модел};
    \node[ormblock, below=of ir] (doc) {Документен\\модел};
    \node[ormblock, below right=of ir] (kv) {Key-value\\модел};
    \node[ormblock, below=of doc, xshift=-3.2cm] (graph) {Графов\\модел};
    \node[ormblock, below=of doc, xshift=3.2cm] (wide) {Wide-column\\модел};

    \draw[ormline] (cpp) -- (ir);
    \draw[ormline] (ir) -- (rel);
    \draw[ormline] (ir) -- (doc);
    \draw[ormline] (ir) -- (kv);
    \draw[ormline] (ir) -- (graph);
    \draw[ormline] (ir) -- (wide);
\end{tikzpicture}
\caption{Таксономия на логическия модел, query IR и backend семействата}
\label{fig:storage_taxonomy}
\end{figure}

Фигура~\ref{fig:storage_taxonomy} подчертава централното теоретично допускане на библиотеката: общото между backend-ите не е техният синтаксис, а логическият модел на данните и операциите. Следователно валидната абстракция трябва да бъде разположена по-високо от SQL, BSON, Redis команди, Cypher или CQL. Тя трябва да работи на равнището на полета, предикати, projection-и, връзки и допустими операции.

\subsection{\ORMName{} като архитектурен подход}

Класическото ORM свързва обектния модел на приложението с неговата персистентна схема~\cite{fowler2002}. В compile-time варианта това свързване се конструира не чрез runtime reflection и конфигурации, а чрез типове, шаблони и \code{constexpr} стойности. Практически това означава, че изборът на полета, допустимите операции и част от съвместимостта с backend-а се проверяват от компилатора. Грешки като обръщение към несъществуващо поле или използване на \code{JOIN} върху конектор без \code{supports\_joins} престават да бъдат runtime изненади.

Съществената особеност е, че compile-time ORM не означава липса на runtime данни. Стойностите на параметрите се подават в момента на изпълнение, но структурата на заявката е вече фиксирана и валидирана. Така се получава ясно разграничение между \emph{логическа форма} и \emph{runtime стойности}. Това разграничение е ключово за тезата, защото позволява едновременно статични гаранции и реална работа с външни бази данни.

От теоретична гледна точка този подход доближава ORM слоя до компилаторен pipeline. Приложният код описва семантична структура, compile-time компонентите я превръщат в междинно представяне, а едва най-долният слой материализира тази структура в backend-специфичен език или протокол. Затова по-късните архитектурни решения в библиотеката могат да се анализират през понятия като междинно представяне, статична проверка и backend materialization.

\subsection{Статичен модел на данните в C++}

В разглежданата библиотека домейн моделът се описва чрез C++ агрегати, чиито полета са от тип \code{property<T, Name>}. По този начин всяко поле носи едновременно стойностен тип и символно име на колоната или атрибута. Допълнително \code{table\_name\_trait<T>} позволява типът да бъде свързан с логически контейнер за съхранение --- таблица, колекция, key prefix, label или друго backend-специфично понятие.

Тази организация има важно теоретично следствие: C++ типът става първичен носител на логическата схема. Отделните backend-и не налагат нов потребителски модел, а само интерпретират един и същ модел по различен начин. Именно затова по-късно може да се говори за еднакъв логически обект, който в релационен backend се превръща в таблица, а в документен --- в документ с вложени или референтни полета.

\begin{table}[H]
\centering
\small
\caption{Съпоставка между C++ конструкциите и теоретичните ORM понятия}
\label{tab:cpp_to_orm_concepts}
\begin{tabularx}{\textwidth}{L{3.3cm}L{4.0cm}Y}
\toprule
\textbf{Конструкция} & \textbf{Теоретична роля} & \textbf{Практически ефект в библиотеката} \\
\midrule
\code{property<T, Name>} & атрибут от логическата схема & носи тип на стойността и символно име за материализация \\
\code{relationship<...>} & логическа връзка между обекти & отделя връзката от физическия механизъм за съхранение \\
\code{table\_name\_trait<T>} & логически контейнер & служи като име на таблица, колекция, key prefix или label \\
\code{field<\&T::m>} & адресиране на атрибут в query пространство & свързва entity слоя с query IR-а \\
\code{connector\_trait<DB>} & backend материализация & определя как логическият модел се изпълнява физически \\
\bottomrule
\end{tabularx}
\end{table}

Следователно ORM слоят не започва от драйвер или от SQL рендериране. Той започва от статично описание на смисъла на данните. Това е и причината entity слоят да бъде фундаментален за цялата архитектура: ако този слой е неясен или непълен, няма как по-долу да се получи коректна и надграждаема материализация.

\begin{lstlisting}[language=C++,caption={Извадка от \code{property} и \code{relationship} в entity слоя},label={lst:property_relationship}]
template <typename T, detail::string_literal Name>
struct property
{
    using value_type = T;

    [[nodiscard]] static constexpr std::string_view column_name() noexcept
    {
        return Name.view();
    }

    T value{};
};

enum class store_as { reference, embed };

template <store_as Strategy, typename T, detail::string_literal Name>
struct relationship
{
    static constexpr store_as strategy = Strategy;
    using related_type = T;

    [[nodiscard]] static constexpr std::string_view field_name() noexcept
    {
        return Name.view();
    }
};
\end{lstlisting}

Listing~\ref{lst:property_relationship} показва, че entity слоят е минималистичен по форма, но достатъчно богат по смисъл. \code{property} не е просто контейнер за стойност; той носи и име, а \code{relationship} не е просто member field; тя фиксира стратегията за материализация като част от типа.

\begin{figure}[H]
\centering
\begin{tikzpicture}[node distance=1.4cm and 1.3cm]
    \node[ormaccent] (entity) {\textbf{Entity тип}\\агрегат от \code{property} и \code{relationship}};
    \node[ormblock, below left=of entity] (props) {Скаларни\\полета};
    \node[ormblock, below right=of entity] (rels) {Логически\\връзки};
    \node[ormblock, below=of props] (schema) {Логическа\\схема};
    \node[ormblock, below=of rels] (physical) {Физическа\\структура};
    \node[ormnote, below=1.2cm of $(schema)!0.5!(physical)$, minimum width=7.8cm] (note)
        {Един и същ C++ модел може да се материализира като SQL таблица,
         документ, key prefix + value, label + properties или wide-column row.};

    \draw[ormline] (entity) -- (props);
    \draw[ormline] (entity) -- (rels);
    \draw[ormline] (props) -- (schema);
    \draw[ormline] (rels) -- (schema);
    \draw[ormline] (schema) -- (physical);
\end{tikzpicture}
\caption{Преминаване от C++ entity описание към логическа схема и физическа материализация}
\label{fig:entity_to_schema_mapping}
\end{figure}

\subsection{Типизиран query IR}

Вместо заявката да се описва директно като SQL или друго текстово представяне, библиотеката използва типизиран query IR. Неговите елементи са \code{field}, \code{Rule}, \code{Placeholder} и query обекти като \code{select\_query}, \code{insert\_query}, \code{update\_query} и \code{delete\_query}. Този модел дава две предимства. Първо, query структурата може да бъде анализирана на compile-time и да участва в \code{static\_assert} проверки. Второ, един и същ IR може да бъде рендиран към различни backend езици и протоколи.

Теоретично IR-ът играе ролята на междинен език между домейн модела и конкретния слой за изпълнение. Това е аналогично на междинните представяния в компилаторите: потребителят работи с високо ниво на абстракция, а backend-ът получава специализирана форма, оптимизирана за собствения си модел. Именно поради това query IR-ът е логически по-близо до семантиката на операциите, отколкото до синтаксиса на конкретна заявка.

\begin{lstlisting}[language=C++,caption={Извадка от \code{select\_query} и композицията на \code{where} клаузи},label={lst:select_query_excerpt}]
template <typename Response, typename Joins, typename Wheres,
          typename Limits, typename Groups, typename Orders>
struct select_query : select_query_tag
{
    using response_type = Response;
    using row_type = projected_type<Response>;

    template <typename RuleType>
        requires is_rule<RuleType>
    [[nodiscard]] constexpr auto where(RuleType r) const
    {
        using NewWheres = detail::append_type_t<Wheres, RuleType>;
        return select_query<Response, Joins, NewWheres, Limits, Groups, Orders>(
            response_, joins_,
            detail::tuple_cat(wheres_, detail::orm_tuple<RuleType>(r)),
            limits_, groups_, orders_);
    }
};
\end{lstlisting}

Listing~\ref{lst:select_query_excerpt} илюстрира важно теоретично свойство: query builder-ът не мутира текст, а изгражда нови типизирани обекти. Това означава, че последователността от операции \code{select(...).where(...).group\_by(...)} може да бъде разглеждана като конструиране на нов тип, а не като поредица от странични ефекти върху низ.

\begin{figure}[H]
\centering
\begin{tikzpicture}[node distance=1.35cm and 1.1cm]
    \node[ormaccent] (query) {\textbf{Потребителска query форма}\\\code{select}/\code{insert}/\code{update}/\code{deleteq}};
    \node[ormblock, below=of query] (field) {\code{field}, \code{Rule}, \code{Placeholder}};
    \node[ormblock, below=of field] (typed) {Типизиран IR\\структура и ограничения};
    \node[ormblock, below left=of typed] (render) {Рендиране\\към backend език};
    \node[ormblock, below right=of typed] (execute) {Изпълнение и\\bind на параметри};

    \draw[ormline] (query) -- (field);
    \draw[ormline] (field) -- (typed);
    \draw[ormline] (typed) -- (render);
    \draw[ormline] (typed) -- (execute);
\end{tikzpicture}
\caption{Pipeline на типизирания query IR от приложния API до backend изпълнението}
\label{fig:query_ir_pipeline}
\end{figure}

\subsection{Теория на връзките и стратегии за съхранение}

Връзките между обекти са един от най-трудните аспекти при обобщаване на ORM извън релационния свят. В релационен контекст една връзка често се материализира чрез външен ключ и \code{JOIN}. В документен контекст обаче същата логическа връзка може да се реализира като вграден документ, а в key-value контекст --- чрез вторичен lookup или денормализиране.

Библиотеката абстрахира тези варианти чрез \code{relationship<Strategy, T, Name>} и \code{store\_as::reference}/\code{store\_as::embed}. \code{reference} обозначава логическа външна препратка, която backend-ът може да реализира чрез foreign key, lookup, отделен ключ или графова връзка. \code{embed} обозначава включване на свързаните данни в същата физическа структура. По този начин връзката се описва веднъж логически, а интерпретацията остава задача на конектора.

\begin{figure}[H]
\centering
\begin{tikzpicture}[node distance=1.5cm and 1.4cm]
    \node[ormaccent] (rel) {\textbf{Логическа връзка}\\entity A $\leftrightarrow$ entity B};
    \node[ormblock, below left=of rel] (reference) {\code{store\_as::reference}\\foreign key / lookup / edge};
    \node[ormblock, below right=of rel] (embed) {\code{store\_as::embed}\\вложен документ / денормализирани данни};
    \node[ormnote, below=3cm of rel, minimum width=8.1cm] (conclusion)
        {Изборът е част от логическото описание, но цената и конкретната реализация
         зависят от backend-а.};

    \draw[ormline] (rel) -- (reference);
    \draw[ormline] (rel) -- (embed);
\end{tikzpicture}
\caption{Двете основни стратегии за материализация на връзки}
\label{fig:relationship_strategies}
\end{figure}

\begin{table}[H]
\centering
\small
\caption{Сравнение между \code{reference} и \code{embed} стратегии}
\label{tab:relationship_strategies}
\begin{tabularx}{\textwidth}{L{2.8cm}L{3.2cm}L{3.0cm}Y}
\toprule
\textbf{Стратегия} & \textbf{Логическа идея} & \textbf{Естествени backend-и} & \textbf{Основен компромис} \\
\midrule
\code{reference} & отделна идентичност на свързания обект & SQL, графови системи, key-based lookup & необходим е допълнителен механизъм за навигация или свързване \\
\code{embed} & физическо включване в родителската структура & документни системи, денормализирани записи & по-трудно независимо обновяване и споделяне на данни \\
\bottomrule
\end{tabularx}
\end{table}

\subsection{Релационен модел}

Релационният модел организира данните в таблици, редове и колони. Основните му предимства са ясна схема, декларативен език за заявки и добре дефинирани ограничения върху целостта. За ORM системите това е естествената среда: свойствата на C++ обекта лесно се съпоставят с колони, а връзките --- с външни ключове. Операции като \code{JOIN}, \code{GROUP BY} и транзакции са част от нормалния езиков модел.

Именно релационният backend е естественият базов сценарий, спрямо който може да се оцени до каква степен една ORM библиотека запазва удобството си и извън SQL света. В настоящата разработка този базов сценарий по-късно се реализира основно чрез \code{SQLiteDB}, а преносимостта на IR-а се анализира допълнително чрез \code{PostgreSQLDB} и \code{MySQLDB}.

\subsection{Документен модел}

Документният модел организира данните в колекции от документи, чиито полета могат да бъдат вложени и не е задължително всички документи да имат еднаква физическа форма. Това го прави естествено съвместим с \code{embed} стратегията и по-гъвкав при денормализирани структури. От друга страна, твърде директното пренасяне на релационни абстракции върху документен backend може да бъде подвеждащо. Не всяка релационна операция има еквивалент със същата цена и семантика.

В контекста на настоящата работа документният модел е важен като тест за това дали query IR-ът описва логика, а не синтаксис на конкретен език. Ако един и същ IR може да се интерпретира като BSON-подобен филтър и projection, без потребителят да пише native document заявки, тогава действително е постигнат storage-agnostic слой на абстракция.

\subsection{Key-value модел}

Key-value системите организират достъпа около ключ и стойност. Те са особено ефективни, когато моделът на достъп е ясен и концентриран около lookup по идентификатор. Недостатъкът е, че общите predicate-базирани заявки и релационните операции са силно ограничени. Затова една ORM абстракция върху key-value backend не може да обещава същата изразителност като върху SQL система.

Този модел е полезен за разграничаване между \emph{логически обект} и \emph{допустим модел на достъп}. Обектът може да бъде един и същ, но не всяка заявка е смислена върху всяка физическа организация. В хранилището това по-късно се вижда ясно при \code{RedisDB}, където lookup логиката е естествена, а join-базираният стил е архитектурно неуместен.

\subsection{Графов модел}

Графовият модел представя данните като възли, ребра и свойства. Силата му е в изразяването на връзки като първокласни обекти, а не като производни от външни ключове. Това го прави подходящ за сценарии със сложни навигации и traversal операции. В същото време графовият модел налага собствен език на заявки и собствена интуиция за селекция и филтриране.

За библиотеката това е показателен сценарий, защото демонстрира, че \code{connector\_trait} може да пренесе общ query IR към \code{MATCH}/\code{RETURN}/\code{WHERE} структура, без потребителят да пише Cypher директно. Ако това е възможно, значи общият слой действително е разположен над конкретния syntax layer.

\subsection{Wide-column модел}

Wide-column системите, каквато е Cassandra, се основават на разпределени таблици с partition и clustering ключове. Макар синтактично да напомнят релационни таблици, те имат различна теория на достъпа: правилната заявка е тази, която следва физическия ключов модел на разпределение. Следователно наличието на поле в C++ модела не означава автоматично, че може да бъде използвано свободно във \code{WHERE} клаузи.

Точно този тип ограничения оправдава capability и constraint подхода на библиотеката. Не е достатъчно query IR-ът да е типово коректен; той трябва да е и допустим спрямо моделa на съхранение. По тази причина wide-column моделът е особено ценен за теоретичната защита на compile-time ограниченията.

\begin{table}[H]
\centering
\small
\caption{Съпоставка между основните модели за съхранение}
\label{tab:storage_models}
\begin{tabularx}{\textwidth}{L{2.3cm}L{2.7cm}L{2.8cm}Y}
\toprule
\textbf{Модел} & \textbf{Единица за съхранение} & \textbf{Типични връзки} & \textbf{Основно ограничение} \\
\midrule
Релационен & таблица/ред & foreign key, join & необходимост от съгласувана схема и силна типова дисциплина \\
Документен & документ & embed или reference & не всяка релационна операция е естествена или евтина \\
Key-value & ключ и стойност & lookup по ключ, hash полета & силно ограничени predicate заявки извън ключовия достъп \\
Графов & възел и ребро & traversal връзки & специфичен query език и графова семантика \\
Wide-column & partitioned row & ключово-ориентирани зависимости & заявките трябва да следват partition логиката \\
\bottomrule
\end{tabularx}
\end{table}

\subsection{Тестова опора на теоретичните твърдения}

Теоретичните твърдения в тази глава не са откъснати от кода. Entity слоят се валидира пряко в \path{tests/unit/test_property.cpp} и \path{tests/unit/test_relationship.cpp}, където се проверяват имената на колоните, разпознаването на entity типове, различаването между \code{store\_as::reference} и \code{store\_as::embed}, както и one-to-many inference логиката. Типизираният query IR е покрит в \path{tests/unit/test_query.cpp}, където се демонстрират field-based предикати, join форми, update composition и delete semantics. Допълнително \path{tests/unit/test_cpp26_reflection.cpp} показва, че compile-time описанието на полетата може да бъде надградено и чрез reflection-aware path, когато toolchain-ът го позволява.

Тази връзка между теория и тестова опора е съществена. Тя показва, че представените понятия не са свободни академични метафори, а реални архитектурни единици, които могат да бъдат проверявани в изходния код и в автоматизираните тестове.

\subsection{Граници на унифицираната абстракция}

След изложеното може да се направи важен извод. Унифицираната ORM абстракция не означава, че всички backend-и стават еднакви. Унифицирането е възможно на равнище логически модел, поле, връзка, query IR и базови операции. Отвъд тази граница започват различията: SQL \code{JOIN} няма универсален смисъл в Redis; embedding логиката няма идентична цена в класическа SQL таблица; graph traversal не се свежда до обикновен \code{SELECT}; partition-driven ограниченията в Cassandra не могат да бъдат маскирани като обикновени predicate правила.

Поради това архитектурно правилният подход не е да се скрият всички различия, а да се направят явни и формално контролирани. Именно тази роля изпълняват \code{connector\_trait}, capability tags и compile-time ограниченията, разгледани в следващите глави. Точно там ще стане видимо как теоретичната рамка от настоящата глава се превръща в работеща архитектура на библиотеката.
